% -----------------------------------------------------------
\section{Key}
% -----------------------------------------------------------
	\label{KEY}
	
% % ----------------------
% \subsection{See also}
% % ----------------------

% ----------------------
\subsection{1828 American Dictionary of the English Language} % *1828
% ----------------------
	\label{KEY-1828}

	\begin{compactitem}
		\item In a general sense, a fastener; that which fastens; as a piece of wood in the frame of a building, or in a chain, etc.
		\item[1] An instrument for shutting or opening a lock, by pushing the bolt one way or the other. Keys are of various forms, and fitted to the wards of the locks to which they belong.
		\item[2] An instrument by which something is screwed or turned; as the key of a watch or other chronometer.
		\item[3] The stone which binds an arch.
		\item[4] In an organ or harpsichord, the key or finger key is a little lever or piece in the fore part by which the instrument is played on by the fingers.
		\item[5] In music, the key or key note, is the fundamental note or tone, to which the whole piece is accommodated, and with which it usually begins and always ends. There are two keys, one of the major, and one of the minor mode. key sometimes signifies a scale or system of intervals.
		\item[6] An index, or that which serves to explain a cypher. Hence,
		\item[7] That which serves to explain any thing difficult to be understood.
		\item[8] In the Romish church, ecclesiastical jurisdiction, or the power of the pope, or the power of excommunicating or absolving.
		\item[9] A ledge or lay of ricks near the surface of the water.
		\item[10] The husk containing the seed of an ash.
	\end{compactitem}

% ----------------------
\subsection{Oxford English Dictionary} % *OED
% ----------------------
	\label{KEY-OED}

	\begin{compactitem}
		\item[II.3.a] With allusion to Matthew 16:19 (see quot. c1384).] \textit{Roman Catholic Church.} The spiritual authority believed to have been transmitted from Christ to St Peter, and so to the Pope considered as his successor, as symbolized by the keys to the kingdom of heaven. Also in a wider sense: the power or authority of the Christian Church or its priests. Usually in \textit{plural}.
		\item[II.3.b] \textit{gen.} A key used to symbolize power, control, or authority, esp. over a particular place; a key as a symbol of office; a position of power or authority.
		\item[II.4.a] Something likened metaphorically to a key in having the power to open or close something else; a thing which provides access or opportunity; a means to a desired objective.
		\item[II.5.a] A means of understanding something unknown, mysterious, or obscure; a solution or explanation.
	\end{compactitem}

% ----------------------
\subsection{Scriptural Usage} % *SCRIPTURES
% ----------------------
	\label{KEY-SCRIPTURES}

	% % -----
	% \subsubsection{Old Testament} % *OT
	% % -----
		% \label{KEY-OT}
	
		% % --
		% \paragraph{KJV}
		% % --
			% \label{KEY-OT-KJV}
	
			% \begin{tabular}{ll}
				% Total occurrences in the OT & 
			% \end{tabular}
			
			% \begin{compactdesc}
				% \item[]
			% \end{compactdesc}
			
		% % --
		% \paragraph{Hebrew Bible}
		% % --
			% \label{KEY-HEBREW}
		
			% %
			% \subparagraph{\texorpdfstring{\he{sample word}}{}}
			% %
				% \label{PAGA} % whatever is in step bible without periods
			
				% \begin{compactdesc}
					% \item[HALOT , PDF ] \textbf{}
						% \begin{compactitem}
							% \item[1]
						% \end{compactitem}
					% \item[BDB , PDF ] \textbf{}
						% \begin{compactitem}
							% \item[1]
						% \end{compactitem}
					% \item[DCH PDF ] \textbf{}
						% \begin{compactitem}
							% \item[1]
						% \end{compactitem}
				% \end{compactdesc}
				
				% \begin{compactdesc}
					% \item[Some OT uses] \textbf{}
						% \begin{compactdesc}
							% \item[]
						% \end{compactdesc}
				% \end{compactdesc}
			
		% % --
		% \paragraph{Septuagint}
		% % --
			% \label{KEY-LXX}
		
			% %
			% \subparagraph{\texorpdfstring{\gr{sample word}}{}}
			% %
		
	% -----
	\subsubsection{New Testament} % *NT
	% -----
		\label{KEY-NT}
	
		% --
		\paragraph{KJV}
		% --
			\label{KEY-NT-KJV}
			
	
			% \begin{tabular}{ll}
				% Total occurrences in the NT & 
			% \end{tabular}
			
			\begin{compactdesc}
				\item[{\hyperref[MATTHEW-16-19]{Matthew 16:19}}]
			\end{compactdesc}
			
		% % --
		% \paragraph{Greek New Testament} % *GREEK
		% % --
			% \label{KEY-GREEK}
		
			% %
			% \subparagraph{\texorpdfstring{\gr{sample word}}{}}
			% %
				% \label{KATALLAGE}
			
				% \begin{compactdesc}
					% \item[BDAG , PDF ] \textbf{}
						% \begin{compactitem}
							% \item 
						% \end{compactitem}
					% \item[CGL , PDF ] \textbf{}
						% \begin{compactitem}
							% \item 
						% \end{compactitem}
					% \item[LSJ , PDF ] \textbf{}
						% \begin{compactitem}
							% \item 
						% \end{compactitem}
				% \end{compactdesc}
				
				% \begin{compactdesc}
					% \item[Some NT uses] \textbf{}
						% \begin{compactdesc}
							% \item[]
						% \end{compactdesc}
				% \end{compactdesc}
		
	% % -----
	% \subsubsection{Book of Mormon} % *BOM
	% % -----
		% \label{KEY-BOM}
	
		% \begin{tabular}{ll}
			% Total occurrences in the BOM & 
		% \end{tabular}
		
		% \begin{compactdesc}
			% \item[]
		% \end{compactdesc}
		
	% -----
	\subsubsection{Doctrine and Covenants} % *DC
	% -----
		\label{KEY-DC}
	
		% \begin{tabular}{ll}
			% Total occurrences in the DC & 
		% \end{tabular}
		
		\begin{compactdesc}
			\item[{\hyperref[DC42-69]{DC 42:69}}]
			\item[{\hyperref[DC90-2]{DC 90:2--4, 6--7}}]
			\item[{\hyperref[DC90-6]{DC 90:6--7}}]
			\item[{\hyperref[DC112-15]{DC 112:15--17}}]
			\item[{\hyperref[DC128-8]{DC 128:8--18}}]
		\end{compactdesc}
		
	% % -----
	% \subsubsection{Pearl of Great Price} % *PGP
	% % -----
		% \label{KEY-PGP}
	
		% \begin{tabular}{ll}
			% Total occurrences in the PGP & 
		% \end{tabular}
		
		% \begin{compactdesc}
			% \item[]
		% \end{compactdesc}
		
% % ----------------------
% \subsection{Key Phrases} % *KEYPHRASES *PHRASES
% % ----------------------
	% \label{KEY-PHRASES}

	% \begin{compactdesc}
		% \item[] \textbf{\#times} | list
			% % \begin{compactdesc}
				% % \item[Bible anchor verses] \phantom{.}
					% % \begin{compactdesc}
						% % \item[Romans 5:5] \gr{}
					% % \end{compactdesc}
				% % \item[Commentary] ---
			% % \end{compactdesc}
	% \end{compactdesc}
	
% % ----------------------
% \subsection{Bible Dictionaries} % *DICTIONARIES
% % ----------------------
	% \label{KEY-BD}

% % ----------------------
% \subsection{Restoration Usage} % *RESTORATION
% % ----------------------
	% \label{KEY-RESTORATION}

	% % -----
	% \subsubsection{Joseph Smith} % *JS
	% % -----
		% \label{KEY-JS}
	
		% \begin{compactdesc}
			% \item[{\hyperref[KEY]{Date/Source}}]
		% \end{compactdesc}
	
	% % -----
	% \subsubsection{Brigham Young} % *BY
	% % -----
		% \label{KEY-BY}
	
		% \begin{compactdesc}
			% \item[{\hyperref[KEY]{Date/Source}}]
		% \end{compactdesc}
	
% ----------------------
\subsection{Commentary and Studies} % *COMMENTARY *STUDIES
% ----------------------
	\label{KEY-COMMENTARY}

	% % -----
	% \subsubsection{Key Points}
	% % -----
		% \label{KEY-KEYS}
	
		% \begin{compactitem}
			% \item
		% \end{compactitem}
		
	% -----
	\subsubsection{The ``keys of hell and of death'' (Revelation 1:18)}
	% -----
		\label{KEY-study}
		
		\begin{compactdesc}
			\item[{\hyperref[LCHRIST-084-L]{Luke 12:5}}] \gr{ὑποδείξω δὲ ὑμῖν τίνα φοβηθῆτε· φοβήθητε τὸν μετὰ τὸ ἀποκτεῖναι ἔχοντα ἐξουσίαν ἐμβαλεῖν εἰς τὴν γέενναν· ναί, λέγω ὑμῖν, τοῦτον φοβήθητε.} \textit{(Some being here possesses \gr{ἐξουσία} to throw someone into hell.)}
			\item[{\hyperref[REVELATION-1-18]{Revelation 1:18}}] \gr{καὶ ὁ ζῶν--- καὶ ἐγενόμην νεκρὸς καὶ ἰδοὺ ζῶν εἰμι εἰς τοὺς αἰῶνας τῶν αἰώνων--- καὶ ἔχω τὰς κλεῖς τοῦ θανάτου καὶ τοῦ ᾅδου.}
		\end{compactdesc}